\chapter{Penentuan Struktur: NMR \texorpdfstring{$^{13}$C}{13C}, MS, dan IR Bersama-sama}\label{ch:b2:structure-determination}

Sebuah laboratorium wewangian menerima sebuah vial berisi cairan tak
berwarna yang berbau melati. Tidak seorang pun menuliskan labelnya. \omterm{def:b2:mass-spec-atomic:spectrum}{Spektrum
massa} memberikan rumus molekulnya, spektrum inframerah gugus-gugus fungsinya,
spektrum karbon-13 kerangkanya, dan spektrum proton cara potongan-potongannya
tersambung; dalam waktu satu jam cairan itu sudah bernama. Bab ini
menambahkan \omterm{def:b2:structure-determination:c13}{NMR karbon-13} pada NMR proton dan spektroskopi inframerah dari
jilid Tahun 1 serta pada \omterm{def:b2:mass-spec-atomic:spectrum}{spektrometri massa} di \cref{ch:b2:mass-spec-atomic},
dan menjadikan keempatnya satu metode.

\begin{recall}
Jilid Tahun 1: pergeseran kimia dan perisai, proton ekuivalen, integrasi,
kopling spin--spin dan aturan $n + 1$, bilangan gelombang inframerah
gugus-gugus fungsi dan daerah sidik jari, derajat ketidakjenuhan,
menyelesaikan struktur dari tiga spektrum. \Cref{ch:b2:mass-spec-atomic}: \omterm{def:b2:mass-spec-atomic:molecular-ion}{ion
molekul}, \omterm{def:b2:mass-spec-atomic:isotope-pattern}{pola isotop}, fragmentasi.
\end{recall}

\section{NMR karbon-13}

Karbon-12, isotop utamanya, tidak mempunyai spin inti dan tidak memberikan
sinyal NMR. Karbon-13 mempunyai spin
$\tfrac12$, seperti proton, tetapi hanya merupakan 1.1~\% karbon alam dan
sinyalnya lebih lemah, inti demi inti, daripada sinyal proton; karena itu
spektrum karbon memerlukan lebih banyak sampel atau lebih banyak pemindaian.
% ledger: iso:C13

\begin{definition}[NMR karbon-13]\label{def:b2:structure-determination:c13}
Yang disebut \emph{NMR karbon-13}\index{NMR karbon-13} merekam resonansi
inti-inti \ce{^{13}C} sebuah sampel. Pengukuran ini biasanya dijalankan
dengan \emph{dekopling pita lebar}\index{dekopling pita lebar}: proton-proton
diiradiasi pada seluruh rentang frekuensinya selama pengukuran, yang
menghilangkan setiap kopling karbon--proton, sehingga setiap jenis karbon
memberikan satu garis tunggal. \emph{Karbon-karbon
ekuivalen}\index{karbon ekuivalen}, yang dipertukarkan oleh sebuah simetri
molekul (atau oleh rotasi yang cepat), memberikan garis yang sama.
\end{definition}

\begin{proposition}[Tanpa kopling karbon--karbon]\label{prop:b2:structure-determination:no-cc-coupling}
Kopling antara atom-atom karbon yang bertetangga tidak muncul dalam spektrum
\ce{^{13}C} pada kelimpahan alami.
\end{proposition}

\begin{proof}
Kopling antara dua karbon hanya terlihat pada molekul-molekul yang kedua
karbonnya adalah \ce{^{13}C}. Untuk dua posisi bertetangga tertentu,
peluangnya $a_{13}^2 = 0.011^2 \approx 1.2 \times 10^{-4}$: sekitar satu
molekul dari delapan ribu. Sinyal setiap karbon hampir seluruhnya berasal
dari molekul-molekul yang tetangganya adalah \ce{^{12}C}, yang tidak terlihat
oleh spektrometer; satelit-satelit yang terkopel seratus kali lebih lemah
daripada derau spektrum biasa.
\end{proof}

\begin{proposition}[Banyaknya garis]\label{prop:b2:structure-determination:signals}
Spektrum \ce{^{13}C} terdekopling memperlihatkan satu garis untuk setiap
himpunan \omterm{def:b2:structure-determination:c13}{karbon ekuivalen}, kecuali jika dua himpunan kebetulan mempunyai
pergeseran yang sama.
\end{proposition}

\begin{proof}[Argumen]
\omterm{def:b2:structure-determination:c13}{Karbon-karbon ekuivalen} berada dalam lingkungan yang identik, sehingga
pergeserannya sama; karbon-karbon yang tidak ekuivalen umumnya berbeda.
Dengan kopling ke proton yang dihilangkan dan kopling antarkarbon yang tidak
ada (\cref{prop:b2:structure-determination:no-cc-coupling}), tidak ada yang
memecah sebuah garis. Tumpang-tindih kebetulan terjadi jika dua karbon yang
berbeda mempunyai pergeseran yang lebih berdekatan daripada resolusinya, hal
yang lazim untuk karbon-karbon cincin benzena di dekat \qty{128}{ppm}.
\end{proof}

Pergeseran tersebar sepanjang sekitar \qty{220}{ppm}, dua puluh kali rentang
proton, sehingga tumpang-tindih lebih jarang daripada dalam spektrum proton,
dan pergeseran itu sendiri sudah menunjukkan jenis karbonnya. Bagan di bawah
memperlihatkan pergeseran yang diukur untuk senyawa-senyawa dalam bab ini.

\begin{omfigure}
\begin{tikzpicture}
\begin{axis}[width=0.86\linewidth, height=4.6cm, xmin=0, xmax=220, ymin=-0.6, ymax=4.6, x dir=reverse,
  xlabel={$\delta$ (ppm)}, label style={font=\small}, tick label style={font=\scriptsize},
  ytick={0,1,2,3,4}, yticklabels={C alkil, C--O, C aromatik, C=O ester, C=O keton},
  xtick={0,20,40,60,80,100,120,140,160,180,200,220}, grid=major, grid style={black!10}]
\addplot[only marks, mark=|, mark size=5pt, thick, omDef] table[x=shift, y=cls] {figdata/out/bachelor-2-nmr-spectra-d.dat};
\end{axis}
\end{tikzpicture}

{\small Pergeseran \ce{^{13}C} terukur butan-2-on, pentan-2-on, etil benzoat,
dan benzil etanoat, dipilah menurut jenis karbon (CDCl$_3$; data PubChem).
Karbonil keton terletak di dekat 209, karbonil ester di dekat
167--171, karbon \omterm{def:b2:huckel:aromatic}{aromatik} di antara 128 dan 137, karbon yang terikat pada
oksigen ester di dekat 61--66, dan karbon alkil di bawah 46.}
% figdata: bachelor-2/nmr-spectra.py part d
% ledger: c13:butanone, c13:pentan-2-one, c13:ethylbenzoate, c13:benzylacetate
\end{omfigure}

\begin{proposition}[Tinggi bukanlah jumlah]\label{prop:b2:structure-determination:intensities}
Dalam spektrum \ce{^{13}C} terdekopling rutin, tinggi garis-garis tidak
sebanding dengan banyaknya karbon yang diwakilinya; karbon tanpa hidrogen
memberikan garis yang lemah.
\end{proposition}

\begin{proof}[Argumen]
Spektrum rutin mengulangi pengukuran lebih cepat daripada relaksasi karbon
yang paling lambat kembali ke kesetimbangan, dan karbon tanpa proton yang
terikat berelaksasi lambat, sehingga sinyalnya sebagian jenuh. Dekopling
juga memperkuat sinyal karbon yang membawa proton, dengan besar yang
bergantung pada lingkungannya. Dalam butan-2-on, garis karbonil adalah yang
paling lemah dari keempat garis walaupun mewakili satu karbon seperti yang
lain. Karena itu integrasi, yang begitu berguna untuk proton, tidak dipakai
untuk spektrum karbon rutin.
\end{proof}

\section{DEPT}

\begin{definition}[DEPT]\label{def:b2:structure-determination:dept}
Yang disebut \emph{DEPT}\index{DEPT} (penguatan tanpa distorsi melalui transfer
polarisasi) adalah serangkaian percobaan karbon-13 yang
memberikan tanda atau ketiadaan pada sinyal-sinyal karbon yang mengikat
proton menurut banyaknya hidrogen. Yang disebut \emph{karbon
kuaterner}\index{karbon kuaterner} tidak membawa hidrogen (terikat pada
empat atom lain, atau karbon karbonil atau \omterm{def:b2:huckel:aromatic}{aromatik} tanpa H); karbon ini
tidak memberikan sinyal DEPT.
\end{definition}

\begin{proposition}[Membaca DEPT]\label{prop:b2:structure-determination:dept}
Dalam DEPT-135, karbon \ce{CH} dan \ce{CH3} mengarah ke atas, karbon
\ce{CH2} mengarah ke bawah, dan \omterm{def:b2:structure-determination:dept}{karbon kuaterner} tidak muncul. Dalam
DEPT-90, hanya karbon \ce{CH} yang muncul.
\end{proposition}

\admitted

Percobaan ini memindahkan magnetisasi dari proton-proton ke karbon tempat
proton itu terikat, melalui serangkaian pulsa frekuensi radio yang sudut
terakhirnya memilih tanggapannya; cara pulsa-pulsa itu bekerja dijelaskan
dalam jilid Tahun 3. Membandingkan spektrum terdekopling dengan DEPT-135 dan
DEPT-90 memilah setiap garis menjadi C, CH, \ce{CH2}, atau \ce{CH3}.

\begin{omfigure}
\pgfplotsset{dept/.style={width=\linewidth, height=2.9cm, ycomb, x dir=reverse, ymin=-1, ymax=1.15,
  ytick=\empty, axis y line=none, axis x line=middle, x axis line style={-}, tick label style={font=\tiny},
  title style={font=\scriptsize, at={(0.0,0.95)}, anchor=north west}, every axis plot/.append style={thick}}}
\begin{minipage}[t]{0.49\linewidth}\centering
{\small butan-2-on}\par
\begin{tikzpicture}
\begin{axis}[dept, xmin=0, xmax=220, xtick={0,50,100,150,200}, title={terdekopling}]
\addplot[omDef] table[x=shift, y=dec] {figdata/out/bachelor-2-nmr-spectra-a.dat};
\end{axis}
\end{tikzpicture}\par
\begin{tikzpicture}
\begin{axis}[dept, xmin=0, xmax=220, xtick={0,50,100,150,200}, title={DEPT-135}]
\addplot[omThm] table[x=shift, y=d135] {figdata/out/bachelor-2-nmr-spectra-a.dat};
\end{axis}
\end{tikzpicture}\par
\begin{tikzpicture}
\begin{axis}[dept, xmin=0, xmax=220, xtick={0,50,100,150,200}, title={DEPT-90}]
\addplot[omMeth] table[x=shift, y=d90] {figdata/out/bachelor-2-nmr-spectra-a.dat};
\end{axis}
\end{tikzpicture}
\end{minipage}\hfill
\begin{minipage}[t]{0.49\linewidth}\centering
{\small etil benzoat}\par
\begin{tikzpicture}
\begin{axis}[dept, xmin=0, xmax=220, xtick={0,50,100,150,200}, title={terdekopling}]
\addplot[omDef] table[x=shift, y=dec] {figdata/out/bachelor-2-nmr-spectra-b.dat};
\end{axis}
\end{tikzpicture}\par
\begin{tikzpicture}
\begin{axis}[dept, xmin=0, xmax=220, xtick={0,50,100,150,200}, title={DEPT-135}]
\addplot[omThm] table[x=shift, y=d135] {figdata/out/bachelor-2-nmr-spectra-b.dat};
\end{axis}
\end{tikzpicture}\par
\begin{tikzpicture}
\begin{axis}[dept, xmin=0, xmax=220, xtick={0,50,100,150,200}, title={DEPT-90}]
\addplot[omMeth] table[x=shift, y=d90] {figdata/out/bachelor-2-nmr-spectra-b.dat};
\end{axis}
\end{tikzpicture}
\end{minipage}

{\small Spektrum \ce{^{13}C} terdekopling (pergeseran dan tinggi terukur,
data PubChem) bersama tanggapan \omterm{def:b2:structure-determination:dept}{DEPT} (tanda dari
\cref{prop:b2:structure-determination:dept}, tinggi skematis). Butan-2-on:
C=O 209.3 (tidak ada dalam \omterm{def:b2:structure-determination:dept}{DEPT}), \ce{CH2} 36.9 (ke bawah), \ce{CH3} 29.4
dan 7.9 (ke atas). Etil benzoat: C=O 166.5 dan karbon cincin yang membawa
ester, 130.6, lenyap; tiga CH \omterm{def:b2:huckel:aromatic}{aromatik} (132.8, 129.6, 128.3) tetap ada
dalam DEPT-90; \ce{OCH2} 60.9 mengarah ke bawah.}
% figdata: bachelor-2/nmr-spectra.py parts a, b
\end{omfigure}

\section{Memadukan teknik-teknik}

Setiap teknik menjawab pertanyaannya sendiri, dan jawaban-jawaban itu
dirakit dalam urutan yang tetap: rumus lebih dahulu, karena rumus membatasi
segala hal lain; lalu gugus-gugus fungsi dan kerangka; lalu
sambungan-sambungannya.

\begin{method}[Memecahkan senyawa yang tidak diketahui]\label{met:b2:structure-determination:solve}
\begin{enumerate}
  \item Rumus: \omterm{def:b2:mass-spec-atomic:molecular-ion}{ion molekul}, $M+1$ untuk banyaknya karbon, \omterm{def:b2:mass-spec-atomic:isotope-pattern}{pola isotop},
        \omterm{def:b2:mass-spec-atomic:nitrogen-rule}{aturan nitrogen}; massa eksak bila tersedia. Derajat
        ketidakjenuhan.
  \item Inframerah: \ce{C=O} (dan jenisnya), \ce{O-H}, \ce{N-H}, \ce{C#N},
        \ce{C-H} \omterm{def:b2:huckel:aromatic}{aromatik} atau alkena di atas \qty{3000}{cm^{-1}}.
  \item Karbon-13 dan \omterm{def:b2:structure-determination:dept}{DEPT}: banyaknya jenis karbon (simetri), kelasnya dari
        pergeserannya, dan C, CH, \ce{CH2}, \ce{CH3} dari \omterm{def:b2:structure-determination:dept}{DEPT}. Bandingkan
        dengan rumusnya.
  \item NMR proton: integrasi, pergeseran, multiplisitas; tetangga-tetangga
        dari aturan $n + 1$.
  \item Rakitlah potongan-potongannya; tuliskan setiap struktur kandidat.
  \item Periksalah setiap kandidat terhadap setiap data, termasuk fragmen
        \omterm{def:b2:mass-spec-atomic:spectrum}{spektrum massa}; pertahankan kandidat yang menjelaskan semuanya.
\end{enumerate}
\end{method}

\begin{omfigure}
\begin{tikzpicture}[x=1cm, y=1cm, st/.style={draw=black!70, rounded corners=2pt, fill=white,
  font=\small, align=center, minimum height=0.9cm, inner sep=3pt, text width=2.3cm}]
\node[st, fill=omDef!10] (ms) at (0,0) {MS\\rumus, $M+1$};
\node[st] (dbe) at (2.9,0) {derajat\\ketidak\-jenuhan};
\node[st, fill=omDef!10] (ir) at (5.8,0) {IR\\gugus fungsi};
\node[st, fill=omDef!10] (c13) at (8.7,0) {\ce{^{13}C}, DEPT\\kerangka};
\node[st, fill=omDef!10] (h1) at (8.7,-1.8) {\ce{^1H}\\sambungan};
\node[st] (cand) at (5.8,-1.8) {struktur\\kandidat};
\node[st, fill=omThm!10] (chk) at (2.9,-1.8) {periksa setiap\\data};
\node[st] (ans) at (0,-1.8) {struktur};
\draw[->, thick] (ms) -- (dbe); \draw[->, thick] (dbe) -- (ir); \draw[->, thick] (ir) -- (c13);
\draw[->, thick] (c13) -- (h1); \draw[->, thick] (h1) -- (cand); \draw[->, thick] (cand) -- (chk);
\draw[->, thick] (chk) -- (ans);
\draw[->, thick, densely dashed] (chk.south) -- ++(0,-0.5) -| node[pos=0.25, below, font=\footnotesize]
  {ada data yang tak terjelaskan: kembali ke kandidat} (cand.south);
\end{tikzpicture}
\par\vspace{4pt}

{\small Urutan kerja untuk senyawa yang tidak diketahui
(\cref{met:b2:structure-determination:solve}). Langkah terakhir adalah yang
paling sering dilewati: struktur yang menyisakan satu puncak tak
terjelaskan belum merupakan jawabannya.}
\end{omfigure}

\section{Kasus-kasus yang dikerjakan}

\begin{example}[Sebuah keton, C$_5$H$_{10}$O]\label{ex:b2:structure-determination:ketone}
Sebuah senyawa yang tidak diketahui mempunyai $M = 86$, pita IR di dekat
\qty{1715}{cm^{-1}}, garis-garis karbon pada 208.9, 45.7, 29.8, 17.4, dan
13.7~ppm, serta sinyal-sinyal proton pada 2.41 (t, 2H), 2.13 (s, 3H), 1.61
(sekstet, 2H), dan 0.91 (t, 3H). Satu derajat ketidakjenuhan, sebuah keton
jenuh (tidak ada proton aldehida di dekat 9.7). Lima karbon, lima garis:
tidak ada simetri. Singlet tiga proton pada 2.13 adalah \ce{CH3} di sebelah
karbonil; urutan triplet, sekstet, triplet adalah gugus propil yang
\ce{CH2}-nya pada 2.41 bersentuhan dengan karbonil. Pentan-2-on,
\ce{CH3COCH2CH2CH3}. Isomernya yang simetris, pentan-3-on, akan
memperlihatkan tiga garis karbon dan tidak ada singlet; \omterm{def:b2:mass-spec-atomic:spectrum}{spektrum massanya}
mempunyai ion McLafferty pada 58 yang tidak dimiliki pentan-3-on
(\cref{ch:b2:mass-spec-atomic}).
% ledger: c13:pentan-2-one, nmr:pentan-2-one, ir:CO-ketone, ms:pentan-2-one
\end{example}

\begin{example}[Sebuah ester aromatik, C$_9$H$_{10}$O$_2$]\label{ex:b2:structure-determination:ester}
Sebuah senyawa yang tidak diketahui dengan rumus \ce{C9H10O2} (lima derajat
ketidakjenuhan) memperlihatkan pita karbonil dan tujuh garis karbon: 166.5
(C), 132.8 (CH), 130.6 (C), 129.6 (CH), 128.3 (CH), 60.9 (\ce{CH2}), dan
14.3 (\ce{CH3}). Spektrum protonnya: 8.05 (m, 2H), 7.35--7.63 (m, 3H), 4.37
(q, 2H), 1.38 (t, 3H). Cincin benzena monosubstitusi (empat garis karbon
\omterm{def:b2:huckel:aromatic}{aromatik}, dua di antaranya masing-masing untuk dua karbon; lima proton
\omterm{def:b2:huckel:aromatic}{aromatik}) menjelaskan empat ketidakjenuhan, karbonil yang kelima. Pasangan
kuartet--triplet adalah gugus etil; \ce{CH2}-nya pada 4.37 (dan pada 60.9
dalam spektrum karbon) terikat pada oksigen. Karbonil ester pada 166.5
terikat pada cincin: kedua proton pada 8.05, di sebelahnya, bergeser ke
medan rendah. Etil benzoat, \ce{C6H5COOCH2CH3}.
% ledger: c13:ethylbenzoate, nmr:ethylbenzoate
\end{example}

\section{Jebakan}

\begin{itemize}
  \item \emph{Proton yang dapat dipertukarkan} (\ce{O-H}, \ce{N-H})
        memberikan sinyal lebar pada pergeseran yang berubah-ubah, sering
        tidak terkopel dengan tetangganya, dan lenyap jika setetes
        \ce{D2O} dikocok bersama sampel.
  \item \emph{Sinyal yang tumpang-tindih}: dua karbon cincin benzena, atau
        beberapa \ce{CH2} sebuah rantai, dapat berbagi satu garis;
        cocokkan banyaknya garis dengan rumus sebelum menyimpulkan simetri.
  \item \emph{Proton diastereotopik}: kedua proton \ce{CH2} di sebelah
        pusat stereogenik tidak ekuivalen, dapat mempunyai pergeseran yang
        berbeda, dan saling terkopel.
  \item \emph{Pola orde kedua}: jika dua proton yang terkopel mempunyai
        pergeseran yang selisihnya kurang dari beberapa kali tetapan
        koplingnya, multiplet-multipletnya saling condong dan aturan
        $n+1$ tidak berlaku; spektrometer bermedan lebih tinggi memulihkan
        pola yang sederhana.
\end{itemize}

\begin{safety}{\ghs[5mm]{GHS02}\,\ghs[5mm]{GHS07}\,\ghs[5mm]{GHS08}}
Butan-2-on mudah terbakar. Magnet spektrometer NMR selalu menyala: tidak
boleh ada perkakas baja, tabung gas, atau kartu magnetik di dekatnya, dan
tidak ada akses bagi orang yang memakai alat pacu jantung.
% ledger: ghs:butanone, ghs:benzylacetate
\end{safety}

\section{Latihan}

\begin{exercise}[$\star$]\label{exo:b2:structure-determination:1}
Berapa garis yang diperlihatkan spektrum \ce{^{13}C} terdekopling setiap
xilena (1,2-, 1,3-, dan
1,4-dimetilbenzena)?
\end{exercise}

\begin{exercise}[$\star$]\label{exo:b2:structure-determination:2}
Ramalkan spektrum DEPT-135 dan DEPT-90 butan-2-ol.
\end{exercise}

\begin{exercise}[$\star$]\label{exo:b2:structure-determination:3}
Garis pada 209, 170, 129, 64, dan 14~ppm paling mungkin milik jenis karbon apa?
\end{exercise}

\begin{exercise}[$\star$]\label{exo:b2:structure-determination:4}
Hitunglah derajat ketidakjenuhan \ce{C8H8O2} dan usulkan struktur yang mengandung cincin benzena.
\end{exercise}

\begin{exercise}[$\star\star$]\label{exo:b2:structure-determination:5}
Sebuah cairan \ce{C4H8O} memperlihatkan pita IR pada \qty{1715}{cm^{-1}} dan
garis-garis karbon pada 209, 37, 29, dan 8~ppm (DEPT-135: 37 ke bawah, 29
dan 8 ke atas, 209 tidak muncul). Identifikasilah cairan itu.
\end{exercise}

\begin{exercise}[$\star\star$]\label{exo:b2:structure-determination:6}
Bagaimana pentan-2-on dibedakan dari pentan-3-on hanya dengan NMR
\ce{^{13}C}? Dengan NMR proton? Dengan \omterm{def:b2:mass-spec-atomic:spectrum}{spektrometri massa}?
\end{exercise}

\begin{exercise}[$\star\star$]\label{exo:b2:structure-determination:7}
Sebuah ester mempunyai $M = 88$ dan sinyal proton pada 4.12 (q, 2H), 2.04
(s, 3H), dan 1.26 (t, 3H) (data latihan). Identifikasilah ester itu.
\end{exercise}

\begin{exercise}[$\star\star$]\label{exo:b2:structure-determination:8}
Bedakan 1,2-diklorobenzena dari 1,4-diklorobenzena dengan NMR \ce{^{13}C}.
\end{exercise}

\begin{exercise}[$\star\star$]\label{exo:b2:structure-determination:9}
Dalam 2-klorobutana, mengapa kedua proton gugus \ce{CH2} tidak ekuivalen?
Apa akibatnya pada spektrum proton?
\end{exercise}

\begin{exercise}[$\star\star\star$]\label{exo:b2:structure-determination:10}
Sebuah senyawa yang tidak diketahui \ce{C9H10O} memperlihatkan pita IR pada
\qty{1690}{cm^{-1}}, garis-garis karbon pada 200.8 (C), 137.0 (C), 132.9
(CH), 128.6 (CH), 128.0 (CH), 31.8 (\ce{CH2}), dan 8.3 (\ce{CH3}), serta
sinyal-sinyal proton pada 7.95 (m, 2H), 7.40--7.55 (m, 3H), 2.98 (q, 2H),
dan 1.22 (t, 3H) (data latihan). Identifikasilah senyawa itu, dan jelaskan
bilangan gelombang karbonilnya yang rendah.
\end{exercise}

\begin{exercise}[$\star\star\star$]\label{exo:b2:structure-determination:11}
Empat isomer \ce{C5H10O2}: asam pentanoat, metil butanoat, etil propanoat,
dan propil etanoat. Berikan untuk masing-masing satu ciri spektrum IR atau
NMR protonnya yang membedakannya dari ketiga isomer lainnya.
\end{exercise}

\begin{exercise}[$\star\star\star$]\label{exo:b2:structure-determination:12}
Sebuah laporan menetapkan garis pada 166.5~ppm etil benzoat sebagai karbon
cincin yang membawa ester dan garis pada 130.6~ppm sebagai karbon karbonil.
Tunjukkan, dengan data \omterm{def:b2:structure-determination:dept}{DEPT} dan bagan pergeseran, bahwa penetapan itu
tidak mungkin benar.
\end{exercise}

\section{Soal: Empat Spektrum Melati}

\begin{problem}[{Soal akhir pekan --- rumus dari ion molekul, pita karbonil
inframerah, karbon-13 dan DEPT, spektrum proton, dan sebuah struktur yang
diperiksa terhadap fragmen-fragmennya}]
\label{pb:b2:structure-determination:1}
Sebuah cairan tak berwarna yang berbau bunga melati memberikan data
berikut. \omterm{def:b2:mass-spec-atomic:spectrum}{Spektrum massa} (NIST
WebBook): $m/z$ 150 (31.7), 151 (3.1), 108 (100), 91 (71.7), 90 (40.6), 43
(37.6). Inframerah (data latihan): pita kuat pada 1740 dan
\qty{1230}{cm^{-1}}, pita lemah tepat di atas \qty{3000}{cm^{-1}}, tidak ada
pita lebar di dekat \qty{3300}{cm^{-1}}. Karbon-13, CDCl$_3$ (data PubChem): 170.70, 136.14, 128.56,
128.24, 66.24, 20.82~ppm. Proton, CDCl$_3$, \qty{90}{MHz}: 7.33 (m, 5H), 5.09 (s, 2H), 2.06 (s, 3H).
% ledger: use:benzylacetate, ms:benzylacetate, c13:benzylacetate, nmr:benzylacetate, ir:CO-ester, iso:C12, iso:C13, mass:C12, mass:H1, mass:O16

\begin{center}
\begin{tikzpicture}
\pgfplotsset{dept/.style={width=0.8\linewidth, height=3.2cm, ycomb, x dir=reverse, ymin=0, ymax=1.15,
  xmin=0, xmax=200, ytick=\empty, axis y line=none, axis x line=bottom, x axis line style={-}, tick label style={font=\tiny},
  title style={font=\scriptsize, at={(0.02,0.82)}, anchor=west}, every axis plot/.append style={thick}}}
\begin{axis}[dept, title={terdekopling}, xtick={0,20,40,60,80,100,120,140,160,180,200}]
\addplot[omDef] table[x=shift, y=dec] {figdata/out/bachelor-2-nmr-spectra-c.dat};
\end{axis}
\end{tikzpicture}
\end{center}
% figdata: bachelor-2/nmr-spectra.py part c

\medskip
\noindent\textbf{Bagian I --- Rumus.}
\begin{enumerate}
  \item Dari rasio 151/150, perkirakan banyaknya karbon.
  \item Apa yang dikatakan oleh \omterm{def:b2:mass-spec-atomic:nitrogen-rule}{aturan nitrogen}?
  \item Dengan sembilan karbon, apa yang tersisa dari massa 150? Usulkan
        rumus dengan oksigen.
  \item Mengapa \ce{C10H14O}, yang juga bermassa nominal 150, kurang mungkin?
  \item Hitunglah derajat ketidakjenuhannya.
  \item Hitunglah \omterm{def:b2:mass-spec-atomic:exact-mass}{massa monoisotopik} rumus yang dipertahankan.
\end{enumerate}

\medskip
\noindent\textbf{Bagian II --- Inframerah.}
\begin{enumerate}[resume]
  \item Tetapkan asal-usul pita pada \qty{1740}{cm^{-1}}. Keluarga senyawa
        manakah yang diisyaratkan oleh posisinya?
  \item Tetapkan asal-usul pita pada \qty{1230}{cm^{-1}}.
  \item Apa yang ditunjukkan oleh pita-pita lemah di atas \qty{3000}{cm^{-1}}?
  \item Apa yang disingkirkan oleh ketiadaan pita lebar di dekat \qty{3300}{cm^{-1}}?
  \item Apakah karbonilnya terkonjugasi dengan cincin? Jelaskan dari bilangan gelombangnya.
\end{enumerate}

\medskip
\noindent\textbf{Bagian III --- Karbon-13 dan \omterm{def:b2:structure-determination:dept}{DEPT}.}
\begin{enumerate}[resume]
  \item Tetapkan setiap garis pada sebuah jenis karbon.
  \item Garis manakah yang mengarah ke bawah dalam DEPT-135?
  \item Garis-garis manakah yang lenyap dalam DEPT-135?
  \item Garis-garis manakah yang tetap ada dalam DEPT-90?
  \item Berapa jenis karbon yang dimiliki cincin benzena monosubstitusi?
        Jelaskan mengapa hanya dua garis CH \omterm{def:b2:huckel:aromatic}{aromatik} yang muncul.
  \item Garis pada 128.24 adalah yang paling tinggi. Apakah garis itu
        mewakili karbon yang paling banyak?
\end{enumerate}

\medskip
\noindent\textbf{Bagian IV --- NMR proton dan struktur.}
\begin{enumerate}[resume]
  \item Tetapkan asal-usul ketiga sinyal proton.
  \item Mengapa sinyal pada 5.09 dan 2.06 berupa singlet?
  \item Mengapa \ce{CH2} pada 5.09 begitu jauh di medan rendah?
  \item Rakitlah strukturnya dan namailah.
  \item Isomernya, metil 2-feniletanoat, \ce{C6H5CH2COOCH3}, mempunyai rumus
        yang sama. Data manakah yang menyingkirkannya?
  \item Jelaskan fragmen-fragmen pada 108, 91, dan 43.
  \item Berapa garis yang akan diperlihatkan oleh spektrum \ce{^{13}C}
        terdekopling yang terpisah sempurna untuk senyawa itu, dan berapa
        di antaranya yang mengarah ke bawah dalam DEPT-135?
\end{enumerate}
\end{problem}
